\section{Compositionality of OGS via \piI}
\label{s:compo}
We now present some compositionality results about \COGS and \AOGS, that can be proved
via the
correspondence between OGS and \piI.

Compositionality of OGS amounts to computing
the set of traces generated by
$\conf{M\subst{x}{V}}$ from the set of traces generated by
$\conf{M}$ and $\conf{\pmap{x}{V}}$.
This is the cornerstone of (denotational) game semantics, where the
combination of $\conf{M}$ and $\conf{\pmap{x}{V}}$
is represented via the so-called `parallel composition plus hiding'.
%
This notion of parallel composition of two processes $P,Q$ plus hiding over a
name $x$ is directly expressible in \piI as the process $\nu x(P | Q)$.
Precisely, 
%\begin{definition}
suppose $\FF,\GG$  are two configurations that agree on their polarity functions,
except on a name $x$;  then
we write 
\vskip .1cm  \mbox{$ $}  \hfill $\nu x(\FF | \GG) $\hfill   $(*)$
\vskip .1cm 
\noindent  for 
the  \piI process $\nu x(\encoConI{\FF} | \encoConI{\GG})$, the 
 \emph{parallel composition plus hiding over $x$}.
%\end{definition}
To define this operation directly at the level of OGS, we would have
to generalize its LTS, allowing internal interactions over a name $x$
used both in input and in output.

% \begin{remark}
% At the level of OGS, we cannot compose directly
%  two configurations such as  
% $\conf{M}$ and $\conf{[x \mapsto V]}$ above. This is because they do not agree on the
% polarity of $x$, since $x$ appear free in $M$ so that it must be an O-name in $\conf{M}$.
%  Generalizing composition to accept such incompatible configurations
% would mean generalizing the OGS LTS, allowing internal interactions that are normally 
% forbidden.
% \end{remark}

As the translation $\cbvSymb$ from OGS configurations into
 $\piI$
 validates the  $\beta_v$ rule, we can prove  that the behaviour of 
 $\conf{M\subst{x}{V}}$
(e.g., its set of traces)
is the same as that of the parallel composition plus hiding over $x$
of $\conf{M}$ and $\conf{\pmap{x}{V}}$.
We use the notation $(*)$  above to express the following two results. 

\begin{theorem} %Removed array to push onto multiple line / have a full line
$
  \encoConI{\conf{\pmap{p}{M\subst{x}{V}},\gamma,\phi\cup\phi'}} 
\ {\bowtie}\
 \nu x(
  \encoConI{\conf{\pmap{p}{M},\emptymap,\phi \uplus \{x\}}}
     | \encoConI{\conf{\gamma \cdot\pmap{x}{V},\phi' \uplus \{x\}}})
$
for $ \bowtie \;  \in \{ \wbOPI, \wbPI, \TEN,\TEPI\}$. 
 \end{theorem}  

\begin{proof}
The theorem is a  consequence of the validity of the $\beta_v$ rule; that
is the property $  \encoConI{ (\lambda x. M) V }  \wbPI   \encoConI{  M \sub V x }  $. 
Indeed, by simple algebraic manipulation one can show 
$$   \encoIa{(\lambda x. M) V} p 
 \wbPI  \nu x( \encoIa M p |   \encoIVa {V} x)$$
which is to say
$$ 
\begin{array}{lc}
  \encoConI{\conf{\pmap{p}{M\subst{x}{V}},\emptymap ,\phi\cup\phi'}} 
 \wbPI
 \nu x(
   \encoConI{\conf{\pmap{p}{M},\emptymap,\phi \uplus \{x\}}} 
     \ \ | \encoConI{\conf{[x \mapsto V],\phi' \uplus \{x\}}})
\end{array} 
$$
The result can then be lifted to a  non-empty environment $ \gamma $ by using the congruence
properties of $\wbPI$, and then to the other relations.
\end{proof}

\begin{corollary}
For any trace $\tr$:
 \begin{enumerate}
 \item  $\conf{M\subst{x}{V},p,\gamma,\phi\cup\phi'} \ArrA{\tr}$
 iff $\nu x(
   \conf{M,p,\emptymap,\phi \uplus \{x\}}
      | \conf{\gamma \cdot\pmap{x}{V},\phi' \uplus \{x\}}) \ArrN{\tr}$;
 \item $\conf{\pmap{p}{M\subst{x}{V}},\gamma,\phi\cup\phi'} \ArrC{\tr}$
 iff $\nu x(
   \conf{\pmap{p}{M},\emptymap,\phi \uplus \{x\}}
      | \conf{\gamma \cdot\pmap{x}{V},\phi' \uplus \{x\}}) \ArrPI{\tr}$.
\end{enumerate}
\end{corollary}

\noindent
Other important properties that we can import in OGS from \piI
are the congruence properties for the \AOGS and \COGS semantics. 
We report the result for  $\TEA$; the same result holds for 
$\wbA,\TEp, \wbC$.
\begin{theorem}
 If $\conf{M} \TEA \conf{N}$
 then for any $\LasV$  context $C$, $\conf{C[M]} \TEA \conf{C[N]}$.
% \rmm{g} = \rmm{a,c}$. 
\end{theorem}

\noindent
The result is obtained  from the congruence properties of \piI, 
Corollary~\ref{c:OGS_COGS}, and the compositionality
of the encoding $\cbvSymb$.

\iflong
The adequacy and congruence result yields soundness (i.e., inclusion) 
of the OGS relations (Alternating and
Concurrent) \wrt contextual equivalence. Soundness of OGS
is usually a non-trivial property, e.g., requiring 
ciu-equivalence~\cite{laird2007}.


